\section{The field and the arithmetic statement}\label{sec:presentation}
Let $n=2m$, $d=n+1$, and put
\[
 S^n=\{x\in\R^d:x\cdot x=1\},\qquad
 P_{\Q}=\{p\in\Q^d:\len p=1\},\qquad
 \rad(p)=p/\sqrt{p\cdot p}.
\]
The tangent plane at $\rad(p)$ consists of the vectors $v$ with
$p\cdot v=0$. We use $\ell^1$ distance for inputs and outputs. A bound
in this norm also bounds the Euclidean norm. Using the radial parameter
keeps the triangulation rational. The vector values stay in the same
ambient coordinates throughout. The classical theorem asks for an
exact zero \cite{wiki}. Approximate-zero search for circuit inputs has
also been studied in computational complexity \cite{gh}.

Here iPRA means first-order intuitionistic arithmetic with a symbol and
defining equations for each fixed primitive recursive (PR) definition,
successor axioms, and quantifier-free induction. A theorem about a
presented field is a schema for fixed PR terms, with free number
parameters allowed. A finite table is itself a number and may be
quantified over. This is also the convention of ProveIt's embedded
object calculus \cite{deriv}. We give the numerical theorem first so the
meaning of a real-valued field is explicit.

\begin{definition}[A field name]\label{def:name}
Fix PR functions $A(p,h)\in\Q^d$ and $\mu(h)\in\N$. The following
rational inequalities are the contract on $p,q\in P_\Q$ and
$h,j\in\N$:
\begin{align}
 j\ge h\quad &\Longrightarrow\quad
       \len{A(p,j)-A(p,h)}\le 2\cdot 2^{-h},                  \tag{C}\\
 \len{p-q}\le2^{-\mu(h)}\quad &\Longrightarrow\quad
       \len{A(p,h)-A(q,h)}\le3\cdot 2^{-h},                   \tag{U}\\
 & |p\cdot A(p,h)|\le2^{-h}.                                 \tag{T}
\end{align}
Encodings are canonical at a rational input, or the evaluator is first
composed with rational reduction. There is no monotonicity assumption
on $\mu$. For each fixed $p$, the sequence $A(p,h)$ is a name of the
field value at $\rad(p)$. Condition (T) permits tangency error in each
rational reading.
\end{definition}

\begin{theorem}[A bounded witness in iPRA]\label{thm:main}
For every field name in dimension $n=2m$ there is a PR function $W$
with the following property provable in iPRA under (C), (U), (T).
Given $k\in\N$, set $t=k+6$, $D=2d\,2^{\mu(t)}$, and $q=W(k)$.
Then
\[
  q\in P_\Q,\quad Dq\in\Z^d,\quad
  \len{A(q,t)}\le\tfrac34\,2^{-k},\qquad
  \forall j\ge t\quad \len{A(q,j)}\le\tfrac{25}{32}\,2^{-k}.
\]
The search evaluates $A$ at a single precision on the signed simplex
grid of denominator $D$. A bound on the number of candidate occurrences is
$2^d\binom{D+d-1}{d-1}$. Repeated points may be scanned or omitted.
\end{theorem}

The last inequality gives a margin for interpreting the answer as a
real vector. Every ordinary tolerance $\varepsilon=u/v>0$ reduces to
a dyadic request, for example $k=v$ for positive integers $u,v$ since
$2^v\ge v$. Section~\ref{sec:round} gives the round-sphere version and
the conversion to a rational unit-vector witness. Sections~\ref{sec:table}
and~\ref{sec:finite-degree} prove the finite obstruction used in the search.
For $n=0$ one may simply take $q=1$, since (T) bounds its only component.

The starting point in ProveIt is the sphere chapter of the general
topology packet \cite{spheres}. Its results S01--S06 establish the
finite antipodal obstruction and a search for exactly tangent,
rational-valued Lipschitz fields. The chapter identifies the
real-valued evaluator as a remaining extension. The proof below
supplies that extension: it corrects a single noisy evaluation
at each proof vertex, permits a general modulus, and does not need a
global bound on the size of the field.
